% Experiment 14: Design and Implementation of Rabin Cryptosystem
\section{Theory}
The Rabin cryptosystem is a public-key encryption scheme whose security is
provably equivalent to the difficulty of integer factorization. Two large
primes $p$ and $q$, both congruent to $3 \bmod 4$ (to simplify square-root
computation), are chosen, and the public key is $n = p \times q$. The
private key is the pair $(p, q)$. Encryption of a message $M$ is extremely
simple:
\[
C = M^2 \bmod n
\]
Decryption computes modular square roots of $C$ modulo $p$ and $q$ separately,
then combines them via the Chinese Remainder Theorem (CRT), yielding four
candidate roots (since squaring is 2-to-1 modulo each prime). Redundancy
added before encryption allows identifying the correct plaintext root.

\section{Algorithm}
\begin{enumerate}
  \item Choose two distinct large primes $p$ and $q$, each congruent to $3 \bmod 4$.
  \item Compute the public key $n = p \times q$.
  \item To encrypt message $M$: compute $C = M^2 \bmod n$.
  \item To decrypt ciphertext $C$: compute $m_p = C^{(p+1)/4} \bmod p$ and $m_q = C^{(q+1)/4} \bmod q$.
  \item Use the Extended Euclidean Algorithm to find $y_p, y_q$ such that $y_p \cdot p + y_q \cdot q = 1$.
  \item Combine using CRT to compute the four candidate roots $r_1, r_2, r_3, r_4$ modulo $n$.
  \item Select the correct root using redundancy/padding information (or by inspection in this demonstration).
  \item Verify that the correct root matches the original plaintext.
\end{enumerate}

\section{Code}
\begin{minted}{cpp}
#include<bits/stdc++.h>
using namespace std;
class rabin {
  private:
    long long p, q, n;
    long long modPow(long long base, long long exp, long long mod) {
      long long result = 1;
      base %= mod;
      while (exp > 0) {
        if (exp & 1) result = (__int128)result * base % mod;
        base = (__int128)base * base % mod;
        exp >>= 1;
      }
      return result;
    }
    bool isPrime(long long num) {
      if (num < 2) return false;
      for (long long i = 2; i * i <= num; i++)
        if (num % i == 0) return false;
      return true;
    }
    long long extGCD(long long a, long long b, long long &x, long long &y) {
      if (b == 0) { x = 1; y = 0; return a; }
      long long x1, y1;
      long long g = extGCD(b, a % b, x1, y1);
      x = y1;
      y = x1 - (a / b) * y1;
      return g;
    }
  public:
    rabin(long long primeP, long long primeQ);
    long long encrypt(long long msg);
    vector<long long> decrypt(long long cipher);
};
rabin::rabin(long long primeP, long long primeQ) {
  p = primeP;
  q = primeQ;
  while (!isPrime(p) || !isPrime(q) || p % 4 != 3 || q % 4 != 3) {
    cout << "p and q must be distinct primes with p, q = 3 (mod 4). Enter again: ";
    cin >> p >> q;
  }
  n = p * q;
  cout << "Public key  n = p*q = " << n << endl;
  cout << "Private key (p, q) = (" << p << ", " << q << ")" << endl;
}
long long rabin::encrypt(long long msg) {
  return (__int128)msg * msg % n;
}
vector<long long> rabin::decrypt(long long cipher) {
  long long mp = modPow(cipher, (p + 1) / 4, p);
  long long mq = modPow(cipher, (q + 1) / 4, q);
  long long yp, yq;
  extGCD(p, q, yp, yq);

  long long r1 = ((yp * p % n) * mq % n + (yq * q % n) * mp % n) % n;
  if (r1 < 0) r1 += n;
  long long r2 = n - r1;

  long long s1 = ((yp * p % n) * mq % n - (yq * q % n) * mp % n) % n;
  if (s1 < 0) s1 += n;
  long long s2 = n - s1;

  return {r1, r2, s1, s2};
}
int main() {
  cout << "Enter two primes p and q, both congruent to 3 mod 4: ";
  long long p, q;
  cin >> p >> q;
  rabin cipher(p, q);

  cout << "\nEnter plaintext message (integer < n): ";
  long long msg;
  cin >> msg;

  long long cipherText = cipher.encrypt(msg);
  cout << "\n--- Encryption ---" << endl;
  cout << "Ciphertext C = M^2 mod n = " << cipherText << endl;

  cout << "\n--- Decryption ---" << endl;
  vector<long long> roots = cipher.decrypt(cipherText);
  cout << "Four possible plaintext roots:" << endl;
  for (int i = 0; i < 4; i++) cout << "  r" << i + 1 << " = " << roots[i] << endl;

  bool found = find(roots.begin(), roots.end(), msg) != roots.end();
  cout << "\nOriginal message was: " << msg << endl;
  return 0;
}
\end{minted}

\section{Result and Discussion - I/O}
\begin{verbatim}
Enter two primes p and q, both congruent to 3 mod 4: 7 11
Public key  n = p*q = 77
Private key (p, q) = (7, 11)

Enter plaintext message (integer < n): 20

--- Encryption ---
Ciphertext C = M^2 mod n = 15

--- Decryption ---
Four possible plaintext roots:
r1 = 64
r2 = 13
r3 = 20
r4 = 57

Original message was: 20
\end{verbatim}

\textbf{Observation:} With $p = 7$ and $q = 11$, both satisfying
$p, q \equiv 3 \pmod 4$, the message $M = 20$ was encrypted to
$C = 15$, and decryption correctly produced four candidate roots, one of
which ($r_3 = 20$) matches the original plaintext. This confirms the
4-to-1 nature of Rabin decryption.

\section{Conclusion}
The Rabin cryptosystem was implemented, illustrating its simple squaring-based
encryption and CRT-based decryption. The main practical challenge is ciphertext
ambiguity: without redundancy, the receiver cannot determine which root is
correct. Despite this, Rabin's security proof (equivalence to integer
factorization) makes it theoretically important.
